%% robustness.tex — Section 5: Robustness and Sensitivity
%%
%% [INSERT] markers requiring filled values:
%%   Table 6: IPW/entropy-balanced reweighted estimates (5.2)
%%   [INSERT: lower bound value], [INSERT: upper bound value] — Lee (2009) bounds (5.3)
%%   [INSERT: whether CI includes zero under worst-case] — Lee (2009) inference (5.3)
%%   Table 7 (or inline): balanced-school-panel estimates (5.4)
%%
%% Once regression output is available, replace all [INSERT: ...] with actual values.

\subsection{Stability Across Specifications}
\label{subsec:stability}

Table~5 presents estimates of Equation~(1) across three progressively richer
specifications. Column~(1) includes only the wave indicator and region fixed
effects. Column~(2) adds student demographic controls (age, gender, rural
residence, and school type). Column~(3) further includes books at home and
SES tertile dummies. The estimated gap between the 2021 and 2018 cohorts
is stable: the raw cross-cohort difference in mathematics is $-0.091$~SD
and the fully controlled estimate in column~(3) is also $-0.091$~SD. For
reading, the corresponding figures are $-0.099$~SD in column~(1) and
$-0.099$~SD in column~(3).

The stability of $\hat{\beta}_1$ across columns implies that observable
differences between the two cohorts account for little of the raw gap.
Students in 2021 were modestly more likely to live in rural areas and
substantially more likely to come from high-SES households (see
Section~\ref{sec:data}), yet controlling for these characteristics barely
moves the estimates. This pattern is consistent with the learning losses
reflecting a common shock that cut across the observable distribution of
students rather than a shift in sample composition alone.

\subsection{Sensitivity to SES Compositional Change}
\label{subsec:ipw}

The most consequential threat to interpreting the cross-cohort gap as a
learning loss is the shift in the SES distribution between waves. The share
of students classified as high-SES increased from 21 percent in 2018 to
35 percent in 2021, a standardized mean difference of 0.340 --- large by
conventional thresholds \citep{austin2011introduction}. Because
higher-SES students tend to score higher on both subjects (the estimated
high-vs.-low SES gap reaches $0.124$~SD in reading; see Table~5), a
sample that is increasingly high-SES in 2021 would mechanically attenuate
any downward shift in scores. The OLS estimates in Table~5 therefore
represent a lower bound on the true learning loss under the standard
monotonicity assumption that SES and test scores are positively related.

To assess the sensitivity of the estimates to this compositional shift, I
reweight the 2021 student observations so that the joint SES distribution
in 2021 matches the 2018 distribution, and then re-estimate Equation~(1)
on the reweighted sample. Weights are constructed using entropy balancing
\citep{hainmueller2012entropy}, which exactly satisfies the first-moment
balance condition on the SES tertile shares. The estimating equation
is unchanged; only the effective sample composition of the 2021 wave
is adjusted.

The reweighted estimates appear in Table~6.%
\footnote{As a check on the entropy balancing approach, I also implement
inverse-probability weighting using a logistic model for the 2021
wave indicator as a function of SES tertile dummies; the results are
virtually identical.}
%% PLACEHOLDER: insert Table 6 (IPW/entropy-balanced estimates) here
The direction of the adjustment is as expected. After reweighting, the
mathematics estimate increases in absolute magnitude to approximately
$-$XX.XX~SD and the reading estimate to %% PLACEHOLDER
$-$XX.XX~SD. %% PLACEHOLDER The increase confirms that
positive SES selection in the 2021 sample compressed the baseline
estimates. The main results in Table~5 should accordingly be read as
conservative: they provide a lower bound on the learning losses that
would be observed under a stable SES composition.

\subsection{Bounding Attrition in Mathematics}
\label{subsec:lee}

Test participation declined from 90.8 percent in 2018 to 85.8 percent
in 2021 for mathematics, a drop of five percentage points. Reading
participation followed a similar pattern. Because the missing students
are not a random subset of the enrolled population, the observed
cross-cohort gap in mean scores could reflect selective non-participation
rather than genuine changes in the distribution of achievement.

I apply the bounding procedure of \citet{lee2009training} to assess the
sensitivity of the mathematics estimates to this differential attrition.
The Lee bounds treat participation as a potential outcome and trim the
2021 sample at the margin of the participation gap to construct
worst-case bounds on the average treatment effect. The lower bound on
the learning loss arises from assuming that the students who dropped
out of testing in 2021 were the highest-achieving students in that
cohort; removing them from the sample makes the 2021 mean artificially
low, maximising the apparent loss. The upper bound assumes the opposite:
the non-participants were the lowest-achieving students, so their
absence makes the 2021 mean artificially high and minimises the
apparent loss.

Under these extreme assumptions, the Lee bounds for mathematics are
$-$XX.XX~SD to $-$XX.XX~SD. %% PLACEHOLDER: lower and upper bounds
XX.XX. %% PLACEHOLDER: one sentence on whether 95% CI excludes zero
For reading, where the participation shortfall is similar in magnitude,
the bounds are $-$XX.XX~SD to $-$XX.XX~SD. %% PLACEHOLDER

Taken together with the SES reweighting results in
Section~\ref{subsec:ipw}, the bounds analysis strengthens the
interpretation that the cross-cohort estimates in Table~5 understate
rather than overstate the true learning loss. Attrition from the 2021
testing round was more likely among lower-performing students, both
because rural schools --- which tend to score lower --- became more
prevalent in the 2021 sample, and because absenteeism is correlated
with socioeconomic disadvantage.

\subsection{Balanced School Panel}
\label{subsec:balanced}

The two assessment waves do not share an identical set of schools. A small
number of schools participated in only one wave: 11 schools appear in
2018 but not in 2021. If school-level characteristics that predict
achievement correlate with the probability of appearing in both waves,
the full-sample estimates could be affected by compositional differences
at the school level as well as the student level.

I address this concern by restricting estimation to the balanced panel of
schools that appear in both waves. Table~7 reports these results.
The point estimates on $\mathit{May2021}_t$ are $-$XX.XX~SD for
mathematics and $-$XX.XX~SD for reading. %% PLACEHOLDER: balanced-panel estimates
Both figures are close to the full-sample estimates in Table~5, indicating
that school-level attrition does not materially alter the cross-cohort
comparison. Standard errors are slightly larger in the balanced-panel
sample, reflecting the reduced sample size, but XX.XX. %% PLACEHOLDER: inference sentence

Across all four robustness exercises --- stable specifications, SES
reweighting, Lee bounds, and the balanced school panel --- the evidence
consistently supports a genuine cross-cohort decline in both mathematics
and reading achievement. The main estimates in Table~5 represent a
conservative characterisation of that decline.
